\section{本卷纲要：一座帝国如何反复接通，又如何断开}

汉从前二〇六年战后分封开始，到二二〇年献帝禅位结束。它不是一条由统一、扩张到衰亡的直线。关中仓储、郡县文书、封国、边郡、商路、宗族和宫门，都曾把分散的人、粮、军队与名分接到一起；它们也都可能在危机中被不同集团分别占有。

本卷先把西汉、新、东汉与汉末放进一条共享年表，再分别考察三个政体。这样安排不是把王莽的九年隔成一段插曲，也不是把二二〇年提前写成三国的开端：两次“受禅”和一次“复汉”都追问同一个问题——诏册和名号怎样同军队、簿籍、道路及地方承认相接。

\begin{sourcenote}[读汉史，先分开年序、宣言和执行]
《史记》《汉书》《后汉书》与《后汉纪》提供本卷的叙事和年序主干，《资治通鉴》帮助连缀前后；它们多由后出的朝廷史家编成。居延、肩水金关、悬泉置等简牍，钱范、碑刻、城址和墓葬可校正某一地点的行政或物质条件，却不能替代全国账簿。下文因此把诏令视为政策语言，把战报视为胜方记录；兵数、人口、边界、族名和人物动机若无独立材料相证，均不写成无疑的事实。
\end{sourcenote}
